\section{La loi générale avant ses réalisations par espèce}
\label{sec:ley-general-masa-accion-autonoma}
\index{masse!loi opératorielle générale}
\index{action!échelle absolue de masse}
\index{tours spectrales!loi globale de masse}

La mécanique dimensionnelle reçoit de HMT deux branches déjà construites. La première
fournit une droite d'action avec deux coordonnées orientées ; la seconde
fournit la généalogie de chaque section persistante. La masse apparaît lorsqu'on
fait agir le caractère généalogique sur une section de la droite de masse.
Cette confluence est antérieure à tout cas particulier et, par conséquent,
l'échelle \(10^{-34}\) appartient à la loi générale et non à une espèce isolée.

\subsection{Branche absolue : du déterminant local à la droite de masse}

Le lecteur déterminantiel local fixe
\begin{equation}
 \Sigma_H=\varphi_{\rm HMT}\alpha_{\rm HMT}^{16},
 \qquad \varepsilon_H:=\left\lfloor\log_{10}\Sigma_H\right\rfloor=-34 .
 \label{eq:ley-masa-sigma-menos34-autonoma}
\end{equation}
La puissance seize provient du conoyau
\(
 \mathbb Z/2\oplus\mathbb Z/2\oplus\mathbb Z/4
\)
du bloc déterminantiel. Sur la droite interne d'action
\(\mathcal L_S\), on obtient les deux sections
\begin{equation}
 \boldsymbol\hbar_{\rm pre}
 =H_5\,10^{-34}\mathcal U_S,
 \qquad
 \boldsymbol\hbar_{\rm ret}
 =\eta_{\rm ret}\,10^{-34}\mathcal U_S,
 \label{eq:ley-masa-hbar-pre-ret-autonoma}
\end{equation}
et les cartes additive et multiplicative de son dédoublement sont
\begin{equation}
 \delta_{2w}=H_5-\eta_{\rm ret},
 \qquad
 R_{\rm act}
 =\frac{\boldsymbol\hbar_{\rm pre}}
        {\boldsymbol\hbar_{\rm ret}}
 =\frac{H_5}{\eta_{\rm ret}}>0 .
 \label{eq:ley-masa-razon-accion-autonoma}
\end{equation}

Soit \(t_0\in\mathcal L_T^+\) la durée élémentaire du calendrier TPK et
\(c_{\rm int}\in\mathcal L_C^+\) sa vitesse interne. La périodicité de
108 itérations de l'endomorphisme temporel CT108 fixe la fréquence angulaire
\(\omega_0=2\pi/(108t_0)\). Les sections absolues antérieure et postérieure au
retour sont donc
\begin{equation}
 \begin{aligned}
 \mathbf m_\star^{\rm pre}
 &=\boldsymbol\hbar_{\rm pre}\,\omega_0c_{\rm int}^{-2},\\
 \mathbf m_\star^{\rm ret}
 &=\boldsymbol\hbar_{\rm ret}\,\omega_0c_{\rm int}^{-2},
 \end{aligned}
 \qquad
 \frac{\mathbf m_\star^{\rm pre}}
      {\mathbf m_\star^{\rm ret}}=R_{\rm act}.
 \label{eq:secciones-absolutas-masa-autonoma}
\end{equation}
Toutes deux appartiennent à
\(
 \mathcal L_M=\mathcal L_S\otimes\mathcal L_T^{-1}
 \otimes\mathcal L_C^{-2}
\).
Ainsi sont séparés avec précision la décade, qui fixe l'échelle absolue, et le
quotient du dédoublement bidirectionnel pré-retour/post-retour, qui
transporte la mémoire de l'évolution close.

Dans les coordonnées internes positives
\(t_0=\widehat t_0\mathcal U_T\) et
\(c_{\rm int}=\widehat c\mathcal U_C\), la réduction scalaire de toute
section dont la réalisation physique est déjà fixée prend la forme
\begin{equation}
 \mathbf m_X^\beta
 =10^{-34}\eta_{\rm ret}\,
  \frac{2\pi}{108\widehat t_0}\,\widehat c^{-2}
  a_X^{(0)}R_{\rm act}^{\kappa_X}\,\mathcal U_M .
 \label{eq:ley-escalar-general-masas-autonoma}
\end{equation}
La coordonnée \(a_X^{(0)}\) provient du caractère et de la hiérarchie harmonique ;
\(\kappa_X\) provient du lecteur de chemin. En particulier, la décade
\(-34\) appartient à la section absolue et n'est pas incorporée dans
\(\kappa_X\).

\subsection{Branche généalogique : registre, caractère et hiérarchie harmonique}

Chaque section survivante parcourt la chaîne
\begin{equation}
 \Omega_{\rm HMT}^{\rm surv}
 \longrightarrow\Gamma_{108}^{\rm enr}
 \longrightarrow\mathcal L_{108}
 \longrightarrow\sigma\in\mathbb Z^5
 \longrightarrow(\sigma,\eta)
 \longrightarrow\chi_{\rm full}(\sigma,\eta)\in\mathbb R_{>0}.
 \label{eq:cadena-genealogica-ley-masa-autonoma}
\end{equation}
Le registre conserve l'ordre des événements ; la signature conserve ses cinq
invariants entiers ; le caractère les évalue sans effacer leur provenance. Son
prolongement spectral vit dans le semi-groupe complet
\begin{equation}
 \mathfrak S=\langle90,120\rangle
 =\{90,120\}\cup\{30r:r\ge6\},
 \qquad
 s_h=q_-^h-q_+^h,
 \quad c_h=q_-^h+q_+^h,
 \label{eq:semigrupo-global-masas-autonoma}
\end{equation}
et prend la forme
\begin{equation}
 \chi_{\rm full}(\sigma,\eta)
 =\chi_0(\sigma)
  \exp\!\left(\sum_{h\in\mathfrak S}\eta_hs_h\right).
 \label{eq:caracter-completo-masas-autonoma}
\end{equation}
Les hauteurs \(90,120,180,210,240,270,360,420\) sont huit témoins
généalogiques particulièrement visibles. La loi appartient à tout
\(\mathfrak S\) ; la récurrence du second ordre de \((s_h,c_h)\) prolonge
les témoins à la hiérarchie harmonique complète.

\subsection{Opérateur universel sur les \(13\) multisections du domaine matériel}

L'atlas fournit \(81\) positions, \(56\) familles de chemins, \(13\)
multisections et \(324\) chemins orientés. Pour une multisection \(F\), soit
\(\mathcal H_F=\bigoplus_{\Gamma\in F}\mathbb C e_\Gamma\). Le caractère
basal et les tours définissent l'opérateur positif
\begin{equation}
 \widehat{\mathcal A}^{(0)}_F e_\Gamma
 =\chi_{\rm full}(\sigma_\Gamma,\eta_\Gamma)e_\Gamma.
 \label{eq:operador-caracter-basal-masas-autonoma}
\end{equation}
Les deux histoires conservées par TPK induisent des lecteurs
\(r\in\{D,\Omega\}\) : \(K_D\) agit sur le registre chronologique de
108 itérations de l'endomorphisme temporel CT108 et \(K_\Omega\), sur le mot
dodécaphasique signé. Si
\(\widehat K_F^r e_\Gamma=\kappa_r(\Gamma)e_\Gamma\), soit
\(F_s\subseteq F\) l'ensemble des chemins de la section ou du bloc physique \(s\), et
définissons
\begin{equation}
 \mathcal H_{F,s}:=\bigoplus_{\Gamma\in F_s}\mathbb C e_\Gamma,
 \qquad
 \widehat{\mathcal A}^{(0)}_{F,s}
 :=\left.\widehat{\mathcal A}^{(0)}_F\right|_{\mathcal H_{F,s}},
 \qquad
 \widehat B_{F,s}^r
 :=\left.\widehat K_F^r\right|_{\mathcal H_{F,s}}.
 \label{eq:restricciones-bloque-ley-masas-autonoma}
\end{equation}
La base diagonale des chemins utilisée par les certificats en vigueur rend invariant
\(\mathcal H_{F,s}\) sous les deux opérateurs et transforme
\(\widehat B_{F,s}^r\) en une restriction autoadjointe. Avec ces types, la loi
générale est
\begin{equation}
 \boxed{
 \widehat{\mathbf M}^{\beta,r}_{F,s}
 =\mathbf m_\star^{\rm ret}\otimes
  R_{\rm act}^{\widehat B_{F,s}^r/2}
  \widehat{\mathcal A}^{(0)}_{F,s}
  R_{\rm act}^{\widehat B_{F,s}^r/2},
 \qquad r\in\{D,\Omega\}.}
 \label{eq:ley-operatorial-general-masas-autonoma}
\end{equation}
L'écriture symétrique préserve la positivité sans présupposer la commutation. Dans la
base simultanément diagonale utilisée par les
certificats en vigueur, elle se réduit à
\begin{equation}
 \widehat{\mathbf M}^{\beta,r}_F
 =\mathbf m_\star^{\rm ret}\otimes
  \widehat{\mathcal A}^{(0)}_F
  \exp\!\left((\log R_{\rm act})\widehat K_F^r\right).
 \label{eq:ley-operatorial-diagonal-masas-autonoma}
\end{equation}
L'équation conserve, en un seul objet, l'échelle absolue, le caractère
global, toute la hiérarchie des tours et le correcteur bidirectionnel dépendant
du chemin. La carte physique de chaque secteur agit ensuite sur le spectre de
\eqref{eq:ley-operatorial-general-masas-autonoma}. Le secteur nul possède son
propre morphisme vers zéro ; les registres composés sont concaténés avant
l'évaluation de la signature ; les secteurs topologiques publient d'abord leurs invariants
de fusion et de tressage.

\begin{theorem}[Covariance et annulation de l'échelle commune]
\label{thm:covariancia-ley-general-masas-autonoma}
La loi~\eqref{eq:ley-operatorial-general-masas-autonoma} est positive et
covariante sous changement unitaire de base dans chaque fibre. Pour deux sections propres
\(X,Y\) réalisées dans la même carte de masse,
\begin{equation}
 \frac{m_Y}{m_X}
 =\frac{\chi_{\rm full}(\sigma_Y,\eta_Y)}
        {\chi_{\rm full}(\sigma_X,\eta_X)}
  R_{\rm act}^{\kappa_r(\Gamma_Y)-\kappa_r(\Gamma_X)}.
 \label{eq:cociente-ley-general-masas-autonoma}
\end{equation}
Les facteurs communs \(10^{-34}\mathcal U_S\), \(2\pi/(108t_0)\) et
\(c_{\rm int}^{-2}\) fixent la section absolue et s'annulent exactement dans
le quotient.
\end{theorem}

\begin{proof}
L'opérateur \(\widehat{\mathcal A}^{(0)}_F\) est positif parce que ses
valeurs propres appartiennent à \(\mathbb R_{>0}\) ; l'exponentielle d'un opérateur
autoadjoint est positive. La conjugaison par une matrice unitaire préserve le
spectre et, par conséquent, la loi. En divisant deux valeurs propres réalisées dans la
même carte, la section commune \(\mathbf m_\star^{\rm ret}\) se simplifie et
il reste~\eqref{eq:cociente-ley-general-masas-autonoma}.
\end{proof}

La fibre électronique centrale est de rang un et constitue le premier
corollaire scalaire de cette loi. Les autres multisections conservent le couple
d'opérateurs, leurs spectres et le morphisme physique qui sélectionne saveur,
couleur, orientation ou composition. Cette hiérarchie permet de publier ensemble le
résultat algébrique, sa réalisation dimensionnelle et la confrontation externe sans
confondre leurs domaines.
